\begin{table}[t]
\centering
\small
\setlength{\tabcolsep}{4pt}
\begin{minipage}[t]{0.45\linewidth}
\centering
\caption{\textbf{Compatibility with Test-Time Scaling.} Accuracy (\%) on GAIA Level~2 with Qwen3.5-4B. \methodname yields consistent gains when combined with existing test-time scaling methods. }
\label{tab:planning_composition}
\vspace{0.5em}
\begin{tabular}{lcc}
    \toprule
    Method & Use Idle? & Accuracy \\
    \midrule
    Sequential Revision  & \xmark & 27.9 \\
    \rowcolor{green!10}
    + \textbf{\methodname}  & \cmark & \textbf{32.2} \\
    \midrule
    Planning & \xmark & 32.2 \\
    \rowcolor{green!10}
    + \textbf{\methodname}  & \cmark  & \textbf{35.3} \\
    \bottomrule
\end{tabular}
\end{minipage}\hfill
\begin{minipage}[t]{0.45\linewidth}
\centering
\caption{\textbf{Effect of Idle-Time Length on Accuracy (\%).} GAIA samples are binned by their ultra-short ratio (the fraction of tool calls shorter than the reasoning step) with Gemini-2.5-Flash. Idle-time availability increases from High to Low.}
\label{tab:ultra-short-impact}
\vspace{1em}
\begin{tabular}{lccc}
    \toprule
    Method & High  & Medium  & Low \\
    \midrule
    Vanilla         & 50.0 & 39.0 & 39.0 \\
    \rowcolor{green!10}
    + \textbf{\mname} & 50.0 & \textbf{45.0} & \textbf{46.0} \\
    \bottomrule
\end{tabular}
\end{minipage}
\end{table}
